\chapter{Target reference}
\label{app:targets}

\section{Every \field{lang} value}

\begin{center}
\small
\begin{longtable}{@{}lL{4.2cm}L{3.4cm}L{2.6cm}@{}}
\toprule
\textbf{\code{lang}} & \textbf{Output} & \textbf{Builds with} &
  \textbf{Overloads} \\
\midrule
\endhead
\lang{python} & pybind11 extension module & off-the-shelf compiler & all \\
\lang{nanobind} & nanobind extension module & off-the-shelf compiler & all \\
\lang{node} & N-API native addon & off-the-shelf compiler & first only \\
\lang{wasm} & Emscripten / embind module & off-the-shelf emsdk & first only \\
\lang{lua} & sol2 module, \code{require}-able & off-the-shelf compiler & first only \\
\lang{julia} & CxxWrap.jl / jlcxx module & compiler + CxxWrap.jl & all \\
\lang{csharp} & C-ABI lib + P/Invoke wrappers & compiler + .NET & first only \\
\lang{java} & C-ABI + FFM wrappers & compiler + JDK & first only \\
\lang{qt} & Qt Widgets inspector & compiler + Qt 6 & all \\
\lang{qml} & QML / QtQuick inspector & compiler + Qt 6 & first only \\
\lang{imgui} & Dear ImGui inspector app & off-the-shelf compiler & --- \\
\lang{rest} & cpp-httplib server + client & \textbf{C++26 toolchain} & first only \\
\lang{openapi} & OpenAPI 3.1 specification & text output & --- \\
\lang{json} & nlohmann (de)serialisation & \textbf{C++26 toolchain} & --- \\
\lang{typescript} & ambient \code{.d.ts} declarations & text output & first only \\
\lang{markdown} & API reference document & text output & --- \\
\lang{html} & styled API reference page & text output & --- \\
\lang{paraview} & Server Manager plugin XML & text output & --- \\
\lang{dynamic} & the IR as runtime data & off-the-shelf compiler & all, at
  call time \\
\bottomrule
\end{longtable}
\end{center}

Each name also accepts a deprecated \code{-expanded} spelling resolving to the
same backend.

\section{Feature support}

Which manifest features reach which backend. A blank cell means the backend
ignores the field rather than failing.

\begin{center}
\small
\begin{longtable}{@{}lccccc@{}}
\toprule
\textbf{\code{lang}} & \field{module\_init} & \field{out\_dir} &
  \field{extensions} & \field{sequences} & \field{out\_params} \\
\midrule
\endhead
\lang{python}     & \yes & \yes & \yes & \yes & \yes \\
\lang{nanobind}   & \yes & \yes & \yes & \yes & \yes \\
\lang{node}       & \yes & \yes & \yes & \yes & \yes \\
\lang{wasm}       & \yes & \yes & \yes & \yes & \yes \\
\lang{lua}        & \yes & \yes &      & \yes & \yes \\
\lang{julia}      & \yes & \yes &      &      &      \\
\lang{csharp}     &      &      &      &      &      \\
\lang{java}       &      &      &      &      &      \\
\lang{qt}         &      &      &      &      &      \\
\lang{qml}        &      &      &      &      &      \\
\lang{imgui}      &      &      &      &      &      \\
\lang{rest}       &      &      &      &      &      \\
\lang{typescript} &      &      & \yes & \yes & \yes \\
\lang{markdown}   &      &      & \yes &      &      \\
\lang{html}       &      &      & \yes &      &      \\
\bottomrule
\end{longtable}
\end{center}

\section{Documentation output}

Where each backend puts the documentation, parameter names and default
arguments harvested from your headers (chapter~\ref{ch:documentation}). A dash
means the target has nowhere to put it, not that the data is missing from the
IR.

\begin{center}
\small
\begin{longtable}{@{}lL{5.6cm}L{5.0cm}@{}}
\toprule
\textbf{\code{lang}} & \textbf{Documentation lands as} &
  \textbf{Names \& defaults} \\
\midrule
\endhead
\lang{python} \lang{nanobind} & docstring, with \code{Args:} /
  \code{Returns:} sections; the class comment is the type's docstring &
  \code{py::arg} / \code{nb::arg} keyword arguments; real Python defaults
  (see the import-time rule) \\
\addlinespace
\lang{typescript} & TSDoc block: prose, \code{@param}, \code{@returns} &
  names in the signature; a defaulted parameter is \code{optional?} \\
\addlinespace
\lang{markdown} \lang{html} & prose, a list of documented arguments, a
  \emph{Returns} line & \code{name: type = default} in the signature \\
\addlinespace
\lang{openapi} & operation \code{description}; per-argument
  \code{description} & argument \code{title}, and \code{default} where the
  value is expressible in JSON \\
\addlinespace
\lang{csharp} & \code{/// <summary>} XML doc & names in the wrapper
  signature (keyword-guarded) \\
\addlinespace
\lang{java} & Javadoc block & names in the wrapper signature
  (keyword-guarded) \\
\addlinespace
\lang{qt} \lang{qml} \lang{imgui} & a field's tooltip & the name labels the
  argument box \\
\addlinespace
\lang{dynamic} & text in the emitted \code{MetaClass} tables & names in the
  parameter tables \\
\addlinespace
\lang{node} \lang{wasm} \lang{lua} \lang{julia} & --- (no docstring slot;
  pair with \lang{typescript}) & --- (the wrappers spell positional locals) \\
\addlinespace
\lang{rest} & --- (the spec is \lang{openapi}'s job) & names in the
  browser client's field placeholders \\
\addlinespace
\lang{json} \lang{paraview} & --- & --- \\
\bottomrule
\end{longtable}
\end{center}

\section{Packaging}

Only two backends emit a \code{make\_wheel.py}, so only those two accept the
per-target \field{wheel} / \field{wheel\_dir} (Section~\ref{sec:wheels}); the
field is an error on any other \field{lang}.

\begin{center}
\small
\begin{tabular}{@{}lL{9.4cm}@{}}
\toprule
\lang{python} & One wheel per CPython version --- pybind11 has no stable-ABI
  mode. Covering several means running \code{make\_wheel.py} once per
  interpreter. \\
\lang{nanobind} & One \code{abi3} wheel covering every CPython 3.12+; below
  3.12, tagged for the exact interpreter that built it. \\
every other & No packaging step. \\
\bottomrule
\end{tabular}
\end{center}

\section{Interop and foreign types}

\begin{center}
\small
\begin{tabular}{@{}lL{9.4cm}@{}}
\toprule
\textbf{\code{lang}} & \textbf{\field{interop}: \code{["eigen"]}} \\
\midrule
\lang{python} & native numpy, through \code{<pybind11/eigen.h>} \\
\lang{nanobind} & native numpy, through \code{<nanobind/eigen/dense.h>}
  (dense only) \\
every other & skips the member --- register the concrete type under
  \field{sequences} / \field{matrices} to give them a flat array \\
\bottomrule
\end{tabular}
\end{center}

\begin{center}
\small
\begin{tabular}{@{}lL{4.4cm}L{4.6cm}@{}}
\toprule
\textbf{\code{lang}} & \textbf{\code{std::filesystem::path}} &
  \textbf{\code{std::shared\_ptr<T>} returned} \\
\midrule
\lang{python} & native (\code{os.PathLike} in, \code{pathlib.Path} out) &
  \code{py::class\_<T, shared\_ptr<T>>} \\
\lang{nanobind} & native & native \\
\lang{wasm} & string, via a copy adapter & \code{.smart\_ptr<\ldots>} on the
  class \\
\lang{lua} & string, via a copy adapter & native (sol2) \\
\lang{node} & string, via a copy adapter & the JS object adopts the pointer;
  untrampolined pointees only \\
\lang{typescript} & \code{string} & the pointee type \\
\bottomrule
\end{tabular}
\end{center}

A \code{shared\_ptr<T>} \emph{parameter} is skipped everywhere --- take
\code{const T\&}.

\section{Build shapes for \code{--build}}

\begin{center}
\small
\begin{tabular}{@{}L{6.0cm}L{7.0cm}@{}}
\toprule
\textbf{Backend} & \textbf{Build} \\
\midrule
\lang{python}, \lang{nanobind}, \lang{rest}, \lang{json}, \lang{lua},
  \lang{imgui} & plain \code{cmake} configure + build \\
\lang{node} & \code{npm install} + \code{npm run build} \\
\lang{wasm} & \code{emcmake cmake} + \code{cmake --build} \\
\lang{julia} & \code{cmake} (locates JlCxx by running \code{julia}) \\
\lang{csharp} & \code{cmake}, then \code{dotnet build} \\
\lang{java} & \code{cmake}, then \code{mvn package} \\
\lang{qt}, \lang{qml} & \code{cmake} with \code{-DQT\_DIR} \\
\lang{markdown}, \lang{html}, \lang{typescript}, \lang{openapi},
  \lang{paraview} & nothing to compile \\
\bottomrule
\end{tabular}
\end{center}

A missing toolchain is a \emph{skip with a note}, not a failure. A missing
\code{dotnet} or \code{mvn} downgrades its backend to ``OK (native only)''.

\section{External requirements}

\begin{center}
\small
\begin{tabular}{@{}lL{9.2cm}@{}}
\toprule
\lang{python} & pybind11 (\code{pip install pybind11}) \\
\lang{nanobind} & nanobind; wheels are \code{abi3} on Python 3.12+ \\
\lang{node} & npm, and \code{cmake-js} \\
\lang{wasm} & an activated emsdk \\
\lang{lua} & Lua 5.1--5.4 or LuaJIT --- \textbf{not} 5.5, which sol2 does not
  support yet. sol2 itself is fetched at configure time \\
\lang{julia} & CxxWrap.jl; the build runs \code{julia} to locate JlCxx \\
\lang{csharp} & a .NET SDK \\
\lang{java} & a JDK with the FFM API \\
\lang{qt}, \lang{qml} & Qt 6 --- but no \code{moc} on the generated code \\
\lang{imgui} & nothing; GLFW and OpenGL3 are auto-fetched \\
\lang{rest} & cpp-httplib and nlohmann/json, fetched on first configure, plus
  the C++26 toolchain \\
\bottomrule
\end{tabular}
\end{center}
